\chapter{Introduction}

%% ITSC Introduction

% Complex agent decisions are often characterized by conflicting objectives.
% An autonomous vehicle, for example, must avoid collisions while also staying
% on the road, reaching a goal position,
% and obeying the speed limit.
% This problem is even
% harder when objectives of multiple agents conflict.
% Agents must determine which strategies are preferable,
% for example, going off-road
% to prevent a crash with another vehicle.
% Such decisions reflect an underlying comparative evaluation
% that ranks possible trajectories according
% to subjective preferences.
% We study the problem
% of comparative evaluations
% in transportation systems 
% from the lens of \emph{preference relations}
% that codify prioritized metrics; metrics
% that human drivers implicitly follow on the road.

% A preference is a ``total subjective comparative evaluation'' \cite{hausman:2011}.
% In the context
% of an autonomous vehicle, a total subjective comparative evaluation is a rule
% that, given two trajectories, decides which
% of the two is preferred from the perspective
% of the controller.
% For multiagent systems too,
% preferences model the multiobjective
% and often conflicting requirements
% that the system must adhere to
% for producing good behavior
% as subjective comparative evaluations (\cref{fig:scenario}).

% \begin{figure}[!t]
% \def\svgwidth{\linewidth}
% \subimport{figures/}{scenario.pdf_tex}
% \caption{Agents operate under shared constraints, but are incentivized by different preferences
% (adapted from Zanardi et al.~\protect\cite{zanardi:2022}, Mendes Filho et al.~\protect\cite{filho:2017}, and Lee et al.~\protect\cite{lee:2025}). On the road, situations arise where a vehicle has to coordinate with other vehicles and must relax a low-priority metric to preserve more important metrics, such as giving way to an ambulance.}
% \label{fig:scenario}
% %\vspace{-1em}
% \end{figure}

% One way to express the tension between prioritized metrics
% of multiple agents is to express them
% as lexicographic optimization problems
% in a game of ordered preference \cite{lee:2025}.
% Games of ordered preference are multi-player games
% with prohibitive computational cost.
% Solving these games with a receding horizon alleviates this cost
% through shorter time horizons and the introduction of feedback.
% However, the receding horizon setting does not address the rapid dimensional increase caused
% by adding players or preference levels.
% Iterated best response~(IBR) efficiently solves multi-player games
% by decomposing them into single-player problems,
% where each player optimizes assuming fixed opponent strategies.
% IBR convergence can be slow due
% to the need to explore large opponent decision spaces.
% We propose ``lexicographic IBR over time,''
% an extension of IBR that incorporates past information,
% available in receding horizon,
% to make predictions about future opponent decisions
% that accelerate convergence
% to generalized Nash equilibria
% for games of ordered preference.

% To study this area, this paper asks:
% \begin{enumerate}[itemsep=0pt,topsep=0pt,label=\textlangle ?\textrangle]
%     \item \emph{%
%         Can we efficiently compute approximate-optimal trajectories
%         for games of ordered preference?}\label{Q1}
% \end{enumerate}

% To make progress on this question, we
% \begin{enumerate}
% \item find that lexicographic IBR over time approximates optimal solutions
% for games of ordered preference; and
% \item design experiments demonstrating that games
% of ordered preference are an effective tool
% for analyzing interactive preferences in transportation systems.
% \end{enumerate}


% RAL

Conflicting objectives are present when having to make a decision. Autonomous agents need to make a choice of what objective should be prioritized. In urban driving scenarios, agents may prefer staying on the lane over staying under the speed limit to avoid hitting pedestrians, while the opposite may be preferred in the countryside.
Games of ordered preference model such scenarios in X way.

Non-cooperative games are often characterized by the lack of information
shared between different agents. In driving scenarios, goal destinations are rarely
broadcasted, enforcing an autonomous agent to make estimations on the
other agents' objectives.

While many works succeed in inferring some of the objectives, such as goal
position, there are no formal methods of estimating the preference order of other agents.

% \begin{figure}[!t]
% \def\svgwidth{\linewidth}
% \subimport{figures/}{scenario.pdf_tex}
% \caption{The order of the game objectives determines the behavior of an
% agent. To this end, estimating that order from partial state 
% observations allows the usage of games of ordered preference in
% environments where there is little to no communications between agents,
% such as traffic scenarios with humans involved.}
% \label{fig:scenario}
% %\vspace{-1em}
% \end{figure}


% To study this area, this paper asks:
% \begin{enumerate}[itemsep=0pt,topsep=0pt,label=\textlangle ?\textrangle]
%     \item \emph{%
%         Can we efficiently infer other agents' preference relations in
%         games of ordered preference?}\label{Q1}
% \end{enumerate}

% To make progress on this question, we
% 1. design a method that estimates preference relations in games of ordered
%    preference; and
% 2. test the previous tool on small-scale robots (cref:fig:scenario), demonstrating
%    the utility of preference estimation.
